La ley de enfriamiento de Newton describe el comportamiento térmico de un cuerpo cuando existe una diferencia de temperatura entre el sistema y el medio que lo rodea. Esta ley establece que la rapidez con la que cambia la temperatura de un objeto es proporcional a la diferencia entre su temperatura instantánea y la temperatura del ambiente. Como consecuencia, el sistema evoluciona espontáneamente hasta alcanzar el equilibrio térmico con el entorno.

Matemáticamente, la ley de enfriamiento de Newton se expresa mediante la ecuación diferencial:

\begin{equation}
\frac{dT}{dt}=-k(T-T_a)
\end{equation}

donde $T$ representa la temperatura del sistema, $T_a$ la temperatura del ambiente y $k$ una constante positiva asociada a las propiedades de transferencia de calor entre el sistema y el medio.

Al resolver esta ecuación diferencial por separación de variables, se obtiene la expresión de la temperatura en función del tiempo:

\begin{equation}
T(t)=T_a+(T_0-T_a)e^{-kt}
\end{equation}

donde $T_0$ es la temperatura inicial del sistema en el instante $t=0$. Esta expresión muestra que la temperatura del cuerpo se aproxima exponencialmente a la temperatura ambiente conforme transcurre el tiempo.

En este informe se estudiará experimentalmente la ley de enfriamiento de Newton, analizando la variación de la temperatura de un sistema con respecto al tiempo y comparando los resultados obtenidos con el modelo teórico.\cite{cengel2002thermodynamics}
